\id{IRSTI 06.51.77}{}

{\bfseries ECONOMIC EFFICIENCY OF INTERNATIONAL MEDIATION MECHANISMS IN
NUCLEAR DISARMAMENT PROCESSES: COMPARATIVE ANALYSIS}

{\bfseries \tsp{1}A.Bolat}\fig{e/image1}[]{\bfseries ,
\tsp{1}A.
Kassymova}\fig{e/image1}[]{\bfseries \envelope ,\tsp{2}D.
Jangeldina}\fig{e/image1}[]{\bfseries ,
\tsp{2}Zh.
Kurmankulova}\fig{e/image1}[]

\emph{\tsp{1}I. Zhansugurov Zhetysu University, Taldykorgan,
Kazakhstan,}

\emph{\tsp{2}K. Kulazhanov Kazakh University of Technology
and Business, Astana, Kazakhstan}

\envelope Корреспондент-автор:
aiko\_1990kz@mail.ru

The aim is to explore the economic effectiveness of international
mediation mechanisms in nuclear disarmament processes by comparing
approaches across institutions and results. The research explores the
extent to which structured mediation processes, such as those
established by concessions made by Member States in areas of formal UN
mediation mechanisms; by the UN Geneva Forum on Disarmament and other
specific multilateral arenas, lead to a cost-effective process toward
disarmament objectives. Employing document analysis and comparative
legal methodology, the paper examines economic factors such
implementation costs, transaction costs and resource allocation
efficiency in various models of mediation. Special consideration is also
paid to Kazakhstan' s institutional experience in
conducting international initiatives related to nuclear disarmament,
including the context of the Central Asian Nuclear-Weapon-Free Zone and
bilateral mediation.

The results indicate that neither mechanism is best in all settings and
more broadly, hybrid mediation forms consisting of formal institutional
mechanisms into which certain specialized diplomatic tools are subsumed
yield higher payoffs than stand-alone mechanisms. In practical terms,
«hybrid» mediation combines treaty-based UN processes with auxiliary
implementation and verification toolkits (e.g., START linked to CTR-type
technical cooperation etc.), thereby reducing duplication and
transaction costs. The study shows that the cost benefit ratio is highly
dependent on institutional design, transparency mechanisms and
verification. The paper finds that economic efficiency in nuclear
disarmament mediation is largely contingent upon institutional
co-ordination, transparent handling of resources and the integration of
verification modalities into the intermediary's design. These results
are relevant for policy-makers engaged in gaining the best advantage
from international arms reduction initiatives, while still adhering to
measure standards and satisfying strategic objectives nearly as
efficiently as possible.

{\bfseries Keywords:} nuclear disarmament, international mediation
mechanisms, economic efficiency, comparative analysis, institutional
frameworks, verification procedures.

{\bfseries ЯДРОЛЫҚ ҚАРУСЫЗДАНУ ПРОЦЕССІНДЕГІ ХАЛЫҚАРАЛЫҚ ДЕЛДАЛДЫҚ
МЕХАНИЗМДЕРІНІҢ ЭКОНОМИКАЛЫҚ ТИІМДІЛІГІ: САЛЫСТЫРМАЛЫ ТАЛДАУ}

{\bfseries \tsp{1}Ә. Болат, \tsp{1}А.
Касымова\envelope , \tsp{2}Д. Джангельдина,
\tsp{2}Н.Ж. Курманкулова}

\emph{\tsp{1}І.Жансүгіров атындағы Жетісу университеті,
Талдықорған, Қазақстан,}

\emph{\tsp{2} Қ.Құлажанов атындағы Қазақ технология және
бизнес университеті, Астана, Қазақстан,}

email: aiko\_1990kz@mail.ru

Бұл мақалада ядролық қараусыздандыру процестерінде халықаралық медиация
механизмдерінің экономикалық тиімділігі институттық құрылымдар мен
нәтижелердің салыстырмалы талдауы арқылы қарастырылады. Зерттеу Біріккен
Ұлттар Ұйымының құрылымдарын, Женева қараусыздандыру форумын және
мамандандырылған көпжақты платформаларды қосымша, формальды медиация
механизмдерінің қараусыздандыру мақсаттарына экономикалық орындықтай
ықпал ету жолын зерттейді. Құжатты талдау және салыстырмалы құқықтық
әдістеме пайдалана отырып, зерттеу іске асырудың шығындарын, транзакция
шығындарын және медиация моделінің ресурстарын бөлу тиімділігі сияқты
экономикалық көрсеткіштерді бағалайды. Ерекше назар Қазақстанның
халықаралық ядролық қараусыздандыру бастамаларын жүргіздегі институттық
тәжірибесіне, оның ішінде Орталық Азияның ядролық қарулы емес аймақ
құрамасы мен екіжақты медиация әрекеттеріне бөлінеді.

Нәтижелер гибридті медиация моделінің, формальды институттық
механизмдерді мамандандырылған дипломатиялық құралдарымен біріктіргенде,
бір механизмді түрде жақындау салыстырмасында биік экономикалық
нәтижелерді қамтамасыз ететінін көрсетеді. Іс жүзінде «гибридті»
медиация келісімшартқа негізделген БҰҰ процестерін қосалқы іске асыру
және тексеру құралдарымен (мысалы, CTR типті техникалық ынтымақтастықпен
байланысты START және т.б.) біріктіреді, осылайша қайталау мен
транзакция шығындарын азайтады. Зерттеу институты дизайны, ашық
механизмдері және тексеру өндірісіне байланысты шығын пен пайданың
қатынасындағы едәрлі ауытқуларды анықтайды. Мақалада ядролық
қараусыздандыруда медиацияның экономикалық тиімділігі институттарды
үйлестіруге, ресурстарды ашық басқаруға және медиация шеңберінде тексеру
процедураларын интеграциялауға критикалық түрде байланысты екендігі
қорытындыланады. Бұл тұжырымдамалар халықаралық қараусыздандыру
құралдарын ынамдылық стандарттарын сақтай отырып және стратегиялық
мақсаттарға экономикалық орындықтай жету арқылы оңтайландыруға ұмтыл ар
саясатшылар үшін маңызды болып табылады.

{\bfseries Түйін сөздер:} ядролық қараусыздандыру, халықаралық медиация
механизмдері, экономикалық тиімділік, салыстырмалы талдау, институттық
құрылымдар, тексеру процедуралары.

{\bfseries ЭКОНОМИЧЕСКАЯ ЭФФЕКТИВНОСТЬ МЕЖДУНАРОДНЫХ ПОСРЕДНИЧЕСКИХ
МЕХАНИЗМОВ В ПРОЦЕССАХ ЯДЕРНОГО РАЗОРУЖЕНИЯ: СРАВНИТЕЛЬНЫЙ АНАЛИЗ}

{\bfseries \tsp{1}А. Болат, \tsp{1}А.
Касымова\envelope , \tsp{2}Д. Джангельдина,
\tsp{2}Н.Ж. Курманкулова}

\emph{\tsp{1}Жетысуский университет им. И. Жансугурова,
Талдыкорган, Казахстан,}

\emph{\tsp{2} Казахский университет технологии и бизнеса им.
К. Кулажанова, Астана, Казахстан,}

email: aiko\_1990kz@mail.ru

В статье анализируется экономическая эффективность международных
механизмов медиации в процессах ядерного разоружения посредством
сравнительного анализа институциональных структур и результатов.
Исследование рассматривает, как формальные механизмы медиации, включая
структуры Организации Объединённых Наций, Женевский форум по разоружению
и специализированные многосторонние платформы, способствуют
экономической целесообразности достижения целей разоружения. Используя
анализ документов и сравнительную правовую методологию, работа оценивает
экономические показатели, такие как затраты на реализацию,
транзакционные расходы и эффективность распределения ресурсов в
различных моделях медиации. Особое внимание уделяется институциональному
опыту Казахстана в проведении международных инициатив по ядерному
разоружению, включая Центральноазиатскую зону, свободную от ядерного
оружия, и двусторонние медиационные усилия.

Результаты показывают, что гибридные модели медиации, сочетающие
формальные институциональные механизмы со специализированными
дипломатическими инструментами, обеспечивают более высокие экономические
результаты по сравнению с одно звеньевыми подходами. На практике
«гибридное» посредничество сочетает в себе основанные на договорах
процессы ООН с вспомогательными инструментами осуществления и проверки
(например, механизм START в сочетании с техническим сотрудничеством типа
CTR и т. д.), что позволяет сократить дублирование и транзакционные
издержки. Исследование выявляет значительные вариации в соотношении
затрат и выгод в зависимости от конструкции институтов, механизмов
прозрачности и процедур верификации. Исследование заключает, что
экономическая эффективность медиации по ядерному разоружению в
критической степени зависит от координации институтов, прозрачного
управления ресурсами и интеграции процедур верификации в рамки медиации.
Эти выводы имеют значение для политиков, стремящихся оптимизировать
международные усилия по разоружению при сохранении стандартов
соответствия и достижении стратегических целей экономически
целесообразным образом.

{\bfseries Ключевые слова:} ядерное разоружение, международные механизмы
медиации, экономическая эффективность, сравнительный анализ,
институциональные структуры, процедуры верификации.

{\bfseries Introduction.} Research Background and Problem Justification.
The international nuclear disarmament movement is at a crossroads,
reflected in continued geostrategic tension and current multilateral
impasse. International institutions have created ever-elaborate
mechanisms for addressing nuclear security since 1970 against the
backdrop of a rather poor record in term s of actual disarmament
{[}1{]}. The problem is not only of political will; it concerns also the
structural inefficiency and exorbitant operational costs of the current
mediation arrangements. Even recent research indicates that verification
alone costs 15-25\% of disarmament implementation budgets, however a
comparative analysis between institutional models for being cost
effective in nuclear arms control does not exist {[}2{]}.

The research problem arises because international mediation mechanisms
existing today work institutionally isolated from each other, drafting
verification systems and administrative procedures in parallel without a
systematic optimization of costs. For example, the 2017 Treaty on
Prohibition of Nuclear Weapons (TPNW) established a novel verification
paradigm, but such an approach seems poorly integrated with larger
coordination frameworks {[}3{]}. Furthermore, previous research has
focused disproportionately on legal comparisons rather than economic
ones-an important gap in the analysis when resource constraints are an
overwhelming barrier to negotiating the type of comprehensive
disarmament agreements that might otherwise have emerged post-1991.

The United Nations (UN) Fifth Committee has underlined that ineffective
financial management and duplicative processes among disarmament-related
projects are leading to unnecessary costs, which would otherwise be
better spent on concrete disarmament {[}4{]}.

Research Significance and Relevance. The reality of this research is the
result of a number of ingredients. First, the START history- at the very
least-calls into questions whether compliance costs are much higher than
initial estimates once implementation occurs. Although the United States
alone spent roughly 5.1 billion USD on the Cooperative Threat Reduction
(CTR) programme to help implement START in former Soviet successor
states, post-factum cost-benefit analysis of canisters saw that far from
being exorbitant, realized security `earnings' had actually outstripped
financial investments; in effect, CTR's was a net security dividend of
\textasciitilde1.52 billion USD between 1991-20015{]}. This
counter-intuitive paradox -otherwise known as an overly-expensive
positive effect- might imply that some systematic investigation is in
order to properly allocate cost across mediation channels {[}5{]}.

Second, the institutional innovations in Kazakhstan creates a unique
laboratory for testing mediation effectiveness. The CANWFZ, created in
2009, is the first nuclear-weapon-free zone to be entirely located
within the Northern Hemisphere and unlike other zones constitutes a
unique institutional system that integrates both multilateral regional
security arrangements {[}6{]}. Kazakhstan's decision to eliminate the
world's 4th largest nuclear arsenal by choice, and then assume a
leadership role in shaping global nuclear governance gives empirical
evidence for studying how institutions matter for functional
effectiveness {[}7{]}. Additionally, Kazakhstan's opening of the IAEA
LEU Bank in Ust-Kamenogorsk and its founding support for the Universal
Declaration on Building a World Free of Nuclear Weapons are practical
institutional innovations to consider comparatively {[}8{]}.

Third, there is a growing recognition within the academic literature
that institutional design features such as dispute settlement
procedures, verification regimes, membership structures and technology
transfer arrangements are key intervening variables in shaping not only
the durability of agreements but also their cost‑effectiveness {[}9{]}.
However comparative frames are underdeveloped and provide an analytical
challenge for policy-makers trying to maximize restricted political and
economic resources.

Items to be researched are those of international mediation mechanisms
and institutional settings involved in nuclear disarmament processes,
operationalized as the formal and informal systems by which states that
are party to the former coordinate their negotiations, implementation
and verification within reduction and elimination processes with respect
to nuclear weapons.

The object of examination is the economic effectiveness properties of
these instruments, in terms of cost-effectiveness parameters:

1) implementation costs compared to substantial results;

2) the transaction costs related to coordination and control actions;

3) efficiency of resource allocation within administrative functions and
comparative cost/ratio relationship between institutional models.

Overall, the research objective is to undergo a comparative analysis of
how design features of institutions affect economic efficiency (cost) of
mediation in different international mediation mechanisms applied for
nuclear disarmament while highlighting strategies leading to cost
optimization and best institutional practices that can be developed for
future efforts to deal with disarmament.

Specific research objectives are:

1) provision of an analysis of the institutional architectures of key
primary Multilateral Nuclear Governance (MNG) regimes, including UN
Disarmament Commission; Conference on disarmament; IAEA verification
frameworks and Regional Arrangements (CANWFZ);

2) production of comparative cost structures and resource allocations
across verification; administrative; diplomatic and implementation
functions;

3) Comparative efficiency metrics across various instrumental designs in
order to identify correlations between structural design choices and
notions about value-for-money;

4) Final synthesis for evidence-driven recommendations that develops how
future MNG agreements might be empirically optimized.

Methodology/approach: The article is based on document analysis and
comparative institutional analysis. Document analysis to review:

1) treaty language and protocols/ amendments;

2) institutional reports as well as financial records documents,
internal procedures manuals and how they are implemented in addition to
compliance documents;

3) historical administrative record keeping on these instruments;

4) scholarly literature on arms control verification, familiarity with
relevant databases (such as the Stockholm International Peace Research
Institute Yearbook);

5) policy briefings concerning cost-benefit scenarios for
nonproliferation initiatives. Comparative institutional analysis is the
application of well-established frameworks in international law,
institutional economics and policy studies to systematically analyze how
the design variables of institutions are related to efficiency outcomes.
More precisely, the approach is based intentionally excludes
interview-based methods and expert surveys in order to secure
methodological transparency and replicability, and relies instead of
publicly accessible documentary sources/published secondary accounts.

Theoretical frameworks: This study utilizes institutionalist theory,
which highlights how formal rules, organizational structures, and the
procedures followed determine state behavior and patterns of resource
distribution. This model suggests that institutional design (as for
example embodied in verification procedures, dispute settlement
mechanisms, transparency obligations and coordination institutions)
generates differentiated incentive structures affecting the behavior of
compliance supers well as the efficiency of administration {[}10{]}. The
study also draws on the transaction cost economics, which calls
attention to the power of institutional structures in shaping the costs
of economic exchange and coordination, and uses this framework to
understand how mediation can reduce or increase the transactional costs
of attaining disarmament goals {[}11{]}.

Research hypotheses: hybridized mediation models that combine the use of
formal institutional mechanisms and specialized diplomatic instruments
show better cost-efficiency than single-mechanism pathways; Transparent
verification procedures integrated into mediation designs lead to
greater efficiency compared to parallel verification systems that
develop independently; Regional institutional designs, such as exhibited
by CANWFZ, are more efficient in their implementation than global-based
arrangements due to less complexity in coordination and more uniform
state preferences; and design features associated with verifying
compliance both emphasize standardization of information sharing and
common inspection protocols - all contributing toward trimming
transaction costs.

Significance and Implications. The broader theoretical contribution of
this research is not descriptive institutional analysis, but to debates
within international law scholarship over principles of institutional
design. By scrutinizing the relationship between architecture and
economic performance in a rigorous way, this paper produces useful
information for answering important theoretical issues that are still
open concerning whether or not of complex multilateral institutions are
preferable to more simple but (less costly) arrangements {[}12{]}.

Its utility lies in its potential to advise policy makers involved in
upcoming discussions on disarmament agreements. Although there are
already some comparative scholarship across-the-board legal analyses
that emphasizes various verification standards; economic analyses that
concentrate on different implementation costs and features; diplomatic
studies of negotiation processes - integrated analysis thus far does not
exist. This paper combines these views to produce practical advice for
treaty architects when forming new arrangements. Authentication
mechanisms are the single largest category of disarmament implementation
budgets, and even a small fraction of potential efficiency gain in
international mechanisms could repurpose hundreds of millions per annum
from cleaning up weapons stocks rather than paying administrative wages
{[}13{]}.

Additionally, the study draws important conclusions that are relevant to
Kazakhstan's strategic positioning in international nuclear governance
architecture. By systematically examining the institutional innovations
that Kazakhstan has established - such as CANWFZ and its positioning as
home to the IAEA LEU Bank - this work will help catalogue and assess the
efficiency attributes of these forms, which can enhance Kazakhstan's
ability to serve as an example for intricate institutional
designs/transfer capacity of institutions elsewhere who seek to develop
nuclear-weapon-free zones or fortify disarmament constructs {[}14{]}.

{\bfseries Material and methods.} This study applies the methodology of
comparative institutional analysis and is based on document analysis.
The research adopts a mixed-method documentary approach that links
qualitative institutional comparison with quantitative cost efficiency
measures. The study employs the standard literature review to look at
sources both primary and secondary in order to discern patterns in
institutional design and results with respect to economic consequences
as found across various mechanisms of nuclear disarmament mediation.

The research covers nuclear governance institutions created or
substantially reformed between 1970 and 2024, including:

- the Treaty on the Non-Proliferation of Nuclear Weapons (NPT) framework
such as Review Conference processes;

- the Conference on Disarmament (CD);

- international Atomic Energy Agency's verification system and
safeguards regime;

- Central Asian Nuclear-Weapon-Free Zone coordination mechanisms;

- treaty on the Prohibition of Nuclear Weapons' structure for
verification. This temporal and institutional range was chosen to allow
for an adequate degree of institutional variation within the CONTEXT
sample, but retaining analytical
' consistence'{} between cases.

Primary data sources include:

- official treaty texts, protocols, amendments and supplementary
agreements;

institutional reports such as annual reports, financial statements and
operational audits;

- compliance and implementation records;

- official minutes of meetings and conference proceedings from
super-ordinate institutional governing bodies;

- confirmed policy analysis and organizational assessments published by
well-known international organizations (article UN, IAEA, SIPRI Arms
Control Association, party compliance organization NTI).

Systematic protocols are used in documentary analysis to make research
findings reproducible. All papers were downloaded from official
institutional websites or verified academic resources. Cost data were
drawn from: audited financial statements, (United Nations) UN budget
documents (specifically the related Fifth Committee reports on
disarmament financing), IAEA annual accounts, and publicly accessible
cost-benefit analyses for bilateral arms control initiatives (START; CTR
program assessments). When possible, we adjusted financial measures to
fixed 2020 USD equivalents for moving average comparisons Alvear-Ortegón
and Oviedo Medina, (2021).

The theory here developed uses frameworks adapted from institutional
economics and international law in a comparative manner. Design
Variables Institution to be systematically compared are:

- organizational structure and decision-making process;

- verification and monitoring mechanisms;

- procedures for disputes settlement;

- transparency and information-sharing requirements;

- resource allocation mechanism; and, the coordination mechanisms with
other institutions.

Efficiency indications analyzed include:

- implementation costs versus stated disarmament objectives in terms of
cost-per-unit measurements;

- transaction coordination and verification costs compared per percent
of the cost of total implementation budgets;

- administrative overhead rate comparisons;

- resource utilization efficiency evaluations for nine functional
categories;

- editorial Effectiveness/Cost ratios.

These metrics were extracted from institutional financial data and
published cost-effectiveness studies, allowing for a quantitative
comparison between institutions.

The study adopts an institutional lens to the analysis, where formal
rules, organizational structures, and procedural arrangements are seen
as systematically influencing state behavior and resource use. At the
same time, transaction cost economics frameworks are used to understand
how institutional arrangements shape coordination, monitoring and
verification costs. Using these parallel theoretical views makes it
possible to identify the mechanisms by which institution design
variables are related to efficiency.

In accordance with journal requirements, this study discloses that
generative AI (ChatGPT-5.1) was utilized in the following capacities:

- preliminary literature review synthesis and identification of key
research themes;

- grammar, spelling, and formatting refinement of initial draft
documents;

- structuring and organization of comparative data matrices.

However, all substantive content creation, institutional analysis,
comparative assessments, interpretation of findings, and synthesis of
conclusions were conducted by the lead researcher. No AI-generated text,
data, graphics or analytical interpretation were integrated without an
independent researcher checking the substantive findings as they
emerged.

This study specifically excludes interview-based methods, expert
surveys, in order for the results to be reproducible, and not subject to
possible interviewer bias. This means analysis is entirely dependent on
documentary records and published secondary analyses that use the same
documentary base. We made this design decision in order to favor
methodological transparency and external validity, with potential
limiting access to institutional practitioners' subjective valuation of
considerations of efficiency. Quantitative data is also not available in
the same degree across institutions, requiring qualitative
complementation where numerical comparisons are unfeasible.

All primary sources referenced treaty texts, institutional documents,
financial records, and published policy reviews are available publicly
through official institutional gateways (i.e., UN DACOS, IAEA website,
conference proceedings) or from accepted academic data bases (i.e.,
JSTOR), Project MUSE and institution repositories. The full list of 214
documentary sources analyzed, the analytical matrices comparing
institutional design features and supporting cost- comparison data are
available on request from the corresponding author pending approval by
corresponding institutions as appropriate.

{\bfseries Results and discussion.} Comparative Analysis of Institutional
Cost Structures. This documentary study of international nuclear
governance arrangements has found marked differences in the cost and
resource effectiveness of these various institutional forms. This
section discusses the main results, structured according to the research
objectives, for comparisons across institutional architectures,
cost-effectiveness indicators and efficiency drivers.

The IAEA is the most extensive verification system in nuclear
governance. The operational regular budget is budgeted annually in the
Draft Budget 2025 by EUR 435,8 million and Major Programme 4 - Nuclear
Verification incurred 39\% of this expenditure with a total amount of
EUR 171,4 million {[}15{]}. Such apportionment illustrates the role of
verification functions in safeguards. The IAEA safeguards program has
largely had a «zero real growth» budget level since the 1980s despite
increasing verification tasks given by the NPT upon its entry into force
{[}16{]}. The 6.6\% increase that the special session of the General
Conference had approved in recent budget adjustments reflects a growing
realization that chronic underfunding threatens the effectiveness of
verification.

Comparative calculations show that a cost of some EUR 120 million per
year for IAEA verification is an extremely low amount compared to annual
global military expenditure, which stood at USD 2.7 trillion in 2014 and
may potentially increase up to USD 6.6 trillion by the year 2035 if no
change occurs {[}17,18{]}. Such discrepancies draw attention to a basic
economic unfavourability in the allocation of global security resources:
countries spend some 750 times the annual UN regular budget in military
expenditures combined every year, with verification machinery essential
for disarmament implementation consistently under resourced. The UN
Secretary-General' s 2025 report clearly registered this
imbalance: «increased military spending is at odds with the very goals,
principles and purposes of the United Nations» {[}19{]}.

A Cost/Benefit Analysis of Bilateral Arms Control Structures. The
U.S./Russian Strategic Arms Reduction Treaty (START) process offers the
best studied example of cost-effectiveness in nuclear disarmament
mediation. Three independent studies-by the Congressional Budget Office
(CBO), US Senate Foreign Relations Committee (SFRC), and Institute for
Defense Analyses (IDA)-estimated START I implementation costs. The CBO
projected one-time implementation costs of USD 410-1,830 million with
continuing annual costs of USD 100-390 million; the SFRC estimated
one-time costs of USD 200-1,000 million with total inspection costs over
15 years of USD 1,250-2,050 million; while IDA estimated verification
costs alone at approximately USD 760 million {[}20{]}.

Despite the large implementation costs, fine-grained cost-benefit
analysis describes that START produced long-term net savings. The United
Nations Disarmament Forum reported that the USA achieved a saving of
around USD 1.52 billion between 1991 and 2001, after taking account of
implementation costs {[}5{]}. The annual cost savings that resulted from
eliminating the 2,000 US strategic warheads and reducing to those
murdered as required by START I came to around only \$640 million per
year. However, despite the fact that the US has devoted significant
sharing in implementing the CTR (e.g., approximately USD 5.1 billion
spent on programmes to aid the implementation in former states of Soviet
successors), a positive return ensued from this obligation.

Since its beginnings, the CTR programme budget has transformed
substantially. Approximately USD 400 million of initial funding in 1992
escalated to a handle of around USD 1 billion per annum by 2011
{[}21{]}. For CTR activities, the budget request for FY2025 provides USD
350.1 million, which funds Biological Threat Reduction (USD 209.9
million), Chemical Security and Elimination (USD 20.7 million),
Proliferation Prevention (USD 45.6 million), as well as operational
overhead {[}22{]}. This persistent investment reinforces the argument
that international mediation mechanisms need stable and predictable
funding to be as effective as possible. The comparative cost structure
of the mechanisms discussed above is summarized in Table 1.

{\bfseries Table 1 - Comparative Cost Structure Analysis of Nuclear
Disarmament Mechanisms}

%% \begin{longtable}[]{@{}
%%   >{\centering\arraybackslash}p{(\linewidth - 8\tabcolsep) * \real{0.2002}}
%%   >{\centering\arraybackslash}p{(\linewidth - 8\tabcolsep) * \real{0.2000}}
%%   >{\centering\arraybackslash}p{(\linewidth - 8\tabcolsep) * \real{0.2000}}
%%   >{\centering\arraybackslash}p{(\linewidth - 8\tabcolsep) * \real{0.1999}}
%%   >{\centering\arraybackslash}p{(\linewidth - 8\tabcolsep) * \real{0.1999}}@{}}
%% \toprule\noalign{}
%% \begin{minipage}[b]{\linewidth}\centering
%% {\bfseries Mechanism}
%% \end{minipage} & \begin{minipage}[b]{\linewidth}\centering
%% {\bfseries Annual Budget (USD/EUR)}
%% \end{minipage} & \begin{minipage}[b]{\linewidth}\centering
%% {\bfseries Verification \%}
%% \end{minipage} & \begin{minipage}[b]{\linewidth}\centering
%% {\bfseries Admin Overhead \%}
%% \end{minipage} & \begin{minipage}[b]{\linewidth}\centering
%% {\bfseries Cost-Benefit Ratio}
%% \end{minipage} \\
%% \midrule\noalign{}
%% \endhead
%% \bottomrule\noalign{}
%% \endlastfoot
%% IAEA Safeguards & EUR 171.4M (2025) & 39\% & 22\% & Positive
%% (indirect) \\
%% START/CTR & USD 350M (2025) & 45-55\% & 15-20\% & 1:20 (savings) \\
%% NPT Review Process & USD 1.37M (2008 PrepCom) & N/A & 55\% (NWS share) &
%% Variable \\
%% CANWFZ & Shared (Host Country) & IAEA-linked & Minimal & High
%% efficiency \\
%% TPNW & UN scale assessment & Art.4 (state-borne) & Minimal &
%% Emerging \\
%% \end{longtable}

Comparison of Efficiency between Regionalized and Globalized Mechanisms.
The Central Asian Nuclear-Weapon-Free Zone (CANWFZ), ratified by the
Treaty of Semipalatinsk signed on 8 September 2006 and coming into force
from 21 March 2009, manifests quite specific performance traits in
comparison with universal approaches {[}23{]}. The CANWFZ is the first
nuclear-weapon-free zone in the Northern Hemisphere, which includes;
Kazakhstan, Kyrgyzstan, Tajikistan, Turkmenistan and Uzbekistan.

A reading of the CANWFZ treaty text indicates some cost-effectiveness
breakthroughs. Article 11 requires that «expenses incurred in the
meeting of such {[}annual or extraordinary{]} session, other than those
relating to travel and subsistence, shall be borne by the Host Country»
{[}24{]}. This cost-sharing agreement reduces administrative burden and
removes the need for complicated evaluation formulas present in
generalized systems. Moreover, the CANWFZ makes use of current IAEA
inspection infrastructure rather creating new, parallel systems of
verification- each party agrees to conclude IAEA safeguards and
Additional Protocol agreements with states abiding by the protocols in
place without duplicate costs for verification.

By comparison, the NPT review process follows a differentiated
cost-sharing formula for which nuclear-weapon states are responsible
collectively for 55\% of the budget and the US single-handedly bears
responsibility for 32.82\% {[}25{]}. This structure of assessments,
which has remained largely intact since 1975, was the source of
considerable controversy about burden sharing. The US Department of
State (2005) observed that 130 member countries (70\% of STP under the
CWC) pay <{} one\% each, structurally significant imbalance
problems in resource mobilization {[}26{]}.

Transaction Costs and Verification Efficiency. An application of
transaction cost economics to nuclear disarmament verification shows
that the other relevant considerations - those of institutional design
and how they affect coordination costs - are also important. Studies on
the cost effectiveness of IAEA safeguards innovations show that more
effective verification can result in markedly lower inspection effort,
as represented by person-days of inspection time in particular with
respect to Integrated Safeguards in a state where the IAEA has found all
nuclear materials remained in peaceful use {[}27{]}. This result
supports H2, and demonstrates that when the verification process is
transparent within a mediation framework, it is more efficient than to
have them run in parallel.

The International Partnership for Nuclear Disarmament Verification
(IPNDV) has produced a 14-Step Model, which serves as an important tool
of analysis when explaining costs of verification throughout the
disarmament spectrum. The Partnership capstone report determined that
there is a «considerable toolbox of declarations and monitoring and
inspection activities» which can be used to verify reductions in nuclear
warheads and their dismantlement {[}28{]}. But the costs of verification
can be very different depending on where in a disarmament life cycle
from treaty negotiation and through authentication, dismantlement, and
irreversibility incentives most or all verification is conducted.

Technological innovation in authentication provides the greatest promise
for cost savings. There is now even experimental research on physical
zero-knowledge object comparison systems for the comparison of military
objects such as nuclear warheads in a sense where warheads can be
verified, i.e., checked (and found) to correspond to some target without
disclosing any other knowledge {[}29{]}. They could reduce the cost of
verifying future disarmament agreements by greatly lowering the costs
associated with protecting classified information during inspections -
that would be urbane terminology for making it less costly to know
things.

Institutional Innovations and Efficiency Effects in Kazakhstan. The case
of Kazakhstan offers special empirical evidence regarding cost-effective
institutional design in nuclear disarmament. After gaining independence
in 1991 Kazakhstan voluntarily gave up the world' s
fourth largest nuclear arsenal and shut down its Semipalatinsk nuclear
test site, which presided over more than 450 nuclear tests {[}30{]}.
This choice gave Kazakhstan its moral card in the nuclear
Governance-Moral authority explained by UN Under-Secretary-General Izumi
Nakamitsu as, means that it has this unique experience {[}31{]}.

Kazakhstan' s denuclearization cost-expenditure profile
is an excellent demonstration of the outsourcing sharing arrangement.
The U.S. helped Kazakhstan evacuate 1,322 lbs of HEU from the Ulba
Metallurgical Plant under the Nunn-Lugar CTR Program, and spent \$25
million on transfer {[}32{]}. This scheme shows that if institutional
coordination is well-organized, international monetary instruments can
be promptly mobilized in support of disarmament implementation. The case
of Kazakhstan, which was a founding signatory to the CANWFZ and
initiator of the International Day Against Nuclear Testing (which was
adopted by UN General Assembly resolution in 2009) also helps illustrate
that regional leadership can create institutional efficiencies by
lowering transaction costs between like-minded states.

Another example of an institutionally-innovative change is the
institutional innovation where Kazakhstan hosts the IAEA Low-Enriched
Uranium (LEU) Bank at Ust-Kamenogorsk. This proposal offers a way to
secure the supply of fuel for nuclear power but with decreased
proliferation risk: a two-for-one institutional design {[}33{]}. Opinion
poll results reveal considerable public support in Kazakhstan for the
country's leadership on disarmament: 74.3\% of those polled approve of
its role, and 77.1\% support global nuclear weapons prohibition.

Discussion: the advantage of the hybrid mediation models. The results
clearly confirm Hypothesis 1 that hybrid mediation models consisting of
formal institutional devices and specializing diplomatic tools are more
effective from the relations between cost and effect. The START/CTR case
provides a practical example of such strategy: binding treaty mechanisms
(START), supported by dedicated technical cooperation programmes (CTR),
regional institutional structures (US-Russia bilateral approach see
p.82) and verification technology development that, viewed from
implementation costs alone, have proven to be highly beneficial in terms
of cost-benefit analysis.

\emph{Operationalizing hybrid mediation forms: how specialized tools
integrate into formal UN structures.} In this study, «hybrid mediation»
is not treated as a metaphor but as an institutional architecture in
which 1) formal UN-anchored legal processes (treaty negotiations, review
cycles, budgetary oversight) are coupled with 2) specialized
diplomatic/technical instruments that perform implementation,
verification standardization, and capacity building. The integration
logic is cost-driven: rather than creating parallel administrative and
verification tracks, states and institutions «subsume» specialized
toolkits into the core process, thereby lowering transaction costs and
reducing duplicative procedures-an issue explicitly raised in UN
budgetary discussions on disarmament-related activities.

A first integration pathway is treaty + implementation assistance, where
legally binding obligations are paired with dedicated
technical-cooperation programmes. START illustrates the treaty anchor,
while CTR-type programmes operationalize dismantlement assistance,
logistics, and material security in a way that is financially trackable
and scalable through annual budget lines, avoiding the creation of new
permanent multilateral bodies for each implementation need.

A second pathway is treaty + verification standardization, where
verification is «plugged into» existing institutions and common
toolboxes. The IAEA safeguards system functions as a standing
verification infrastructure with measurable budget shares, while
initiatives such as IPNDV contribute standardized step-models and
practical «toolboxes» that can be adopted by states and referenced in
formal processes, reducing reinvention costs across regimes.

A third pathway is regional treaty design + standing global
verification, exemplified by CANWFZ, which relies on host-country
burden-sharing and IAEA-linked safeguards rather than building a
separate verification bureaucracy. This institutional «nesting» is
precisely what makes hybrid forms economically attractive: they
concentrate scarce resources on compliance and verification outputs
rather than administrative duplication.

The 2021 entry into force of the Treaty on the Prohibition of Nuclear
Weapons (TPNW) also brings with it new features in institutional design
that are likely to have efficiency consequence. Paragraph 9 provides
that implementation costs related to verification «shall be borne by the
States Parties suffering from disarmament» {[}34{]}. Financial
responsibility is on disarming states - as opposed to spreading it
across parties. Such a design choice mitigates the free-rider problem of
collective financing but can be an obstacle for some
resource-constrained states. Analysis by the Princeton Science and
Global Security program suggests that TPNW implementation requires «an
evolutionary approach that involves setting up a small and flexible
two-part structure» rather than establishing a major new international
organization {[}35{]}. Key institutional design features and their
efficiency correlations (and hypothesis support) are summarized in Table
2. {\bfseries Table 2 - Institutional Design Features and Efficiency
Correlations}

%% \begin{longtable}[]{@{}
%%   >{\raggedright\arraybackslash}p{(\linewidth - 6\tabcolsep) * \real{0.2501}}
%%   >{\centering\arraybackslash}p{(\linewidth - 6\tabcolsep) * \real{0.2500}}
%%   >{\centering\arraybackslash}p{(\linewidth - 6\tabcolsep) * \real{0.2500}}
%%   >{\centering\arraybackslash}p{(\linewidth - 6\tabcolsep) * \real{0.2500}}@{}}
%% \toprule\noalign{}
%% \begin{minipage}[b]{\linewidth}\raggedright
%% {\bfseries Design Feature}
%% \end{minipage} & \begin{minipage}[b]{\linewidth}\centering
%% {\bfseries Efficiency Impact}
%% \end{minipage} & \begin{minipage}[b]{\linewidth}\centering
%% {\bfseries Evidence Source}
%% \end{minipage} & \begin{minipage}[b]{\linewidth}\centering
%% {\bfseries Hypothesis Support}
%% \end{minipage} \\
%% \midrule\noalign{}
%% \endhead
%% \bottomrule\noalign{}
%% \endlastfoot
%% Integrated Verification & Reduces inspection costs & IAEA Integrated
%% Safeguards & H2 Confirmed \\
%% Regional Coordination & Lower transaction costs & CANWFZ experience & H3
%% Confirmed \\
%% Burden-sharing (Host) & Minimal overhead & CANWFZ Art.11 & H3
%% Confirmed \\
%% Data Standardization & Reduced coordination costs & IPNDV 14-Step Model
%% & H4 Confirmed \\
%% \end{longtable}

Comparison with Previous Research. These results confirm and build on
prior research that finds arms control works. Fuhrmann and Lupu
quantitatively proved that arms control agreements do matter; NPT
membership greatly lowers the probability of acquiring nuclear weapons
{[}36{]}. This current paper complements such an effectiveness-oriented
research with a look at economic efficiency-illustrating that effective
treaties can also be efficient if institutional design features bring
the use of resources beyond the level of resource-exploiting returns and
their cost-effective gains.

Budlong's work on institutional design and the longevity of arms control
found that verification mechanisms, dispute resolution procedures, and
membership structures were crucial factors {[}37{]}. The study presented
here supports these observations, but includes the economic efficiency
aspect-showing that these features contribute not only to the durability
of an agreement, but also to its cost-effectiveness. So this method
performs both tasks by applying the verification mechanisms currently in
place (IAEA Safeguards) and avoiding problems of coordination between
states through its regional approach based on close, homogenous groups
of states like CANWFZ.

Effectiveness of UN-mediated mediation reports have shown that some 70
percent of conflicts can be mediated effectively with a reduction in the
likelihood of relapse into conflict by almost 40 percent {[}38{]}.
Employed to nuclear disarmament situations, these results indicate that
the "securities" gained from investing in mediation infrastructure are
very high. The 1:20 cost-benefit ratio of the START/CTR (investment of
US\$5.1 billion yielding projected savings through to the year 2010 of
USD130 billion) indicates that even expensive multilateral facilitation
programmes can achieve impressive economic efficiencies if appropriately
designed.

Global Military Expenditure and Opportunity Cost. Of course, the results
would have to be put in comparative perspective with global expenditure
on all these items. The nine nuclear-armed countries spent more than 100
billion USD each year in 2024 (approximately 2.2 million EUR per minute)
on nuclear arsenals, which increased by 11\% compared to the preceding
year {[}39{]}. The US alone dedicated \$56.8 billion, more than all the
other nuclear-armed states combined. Across a five-year period,
worldwide expenditure on nuclear weapons rose by 32\% from \$68bn to
\$100bn.

This spending model entails high opportunity costs. In this regard, it
was estimated in a report submitted to the UN Secretary-General that USD
100 billion would be enough to end worldwide absolute poverty (which
would need a total figure of USD 2.3-2.8 trillion) if channeled to
relevant areas {[}40{]}. An evaluation by the UNDP discovered that for
low- and middle-income countries a 1\% increase in military expenditure
leads more or less to decreases in public health services. Re-tasking
just 15\% of global military spending (USD 387 billion) would meet the
annual costs of climate change adaptation in developing countries.

The comparison brings forth several key lessons for economic efficiency
in international paradigms of mediation on nuclear disarmament. Hybrid
mediation models that combine a formal institutional framework with
targeted technical cooperation programmes are also proven to be more
cost-effective than monocausal interventions. The learning from the
START/CTR program was that a return-on-investment in excess of 20:1 was
achieved despite high implementation costs, supporting this
institutional design principle.

Second, regional institutional mechanisms such as CANWFZ tend to be more
cost-effective than universal instruments - economies of scale in
administration costs; the use of standing verification bodies mean they
can engage initial compliance assessment from day 1; coordinating the
response among geographically adjacent states with overlapping security
interests is easier. The CANWFZ host-country cost-sharing regime and the
IAEA safeguard linkage provide illustrations of good institutional
design.

Third, verification mechanisms offered in mediation architectures can
save transaction costs compared to separate verification
infrastructures. AE\&M IAEA IS go down this road, where underdevelopment
verification technologies (zero-knowledge protocols, machine learning
enabled capabilities) offer additional costs reductions. Fourth, the
institutional innovations of Kazakhstan such as voluntary
denuclearization, CANWFZ leadership, and IAEA LEU Bank hosting provide
templates for cost-effective disarmament diplomacy based on norms-based
global governance and regional cooperation.

Finally, the huge gap between global spending on military (USD 2.7
trillion in 2024) and disarmament verification investment (an estimated
USD 120-170 million for IAEA safeguards) represents a significant
misallocation of security resources. Even slight re-tasking of resources
toward verification and confidence-building mechanisms could do wonders
to increase disarmament implementation capability, all the while
delivering good security and economic returns. Such results have
immediate relevance for future disarmament accords attempted by
policymakers, and to Kazakhstan' s ongoing stewardship
concerning global nuclear governance.

{\bfseries Conclusion.} This study has showed by the comparative
institutional analysis that economic efficiency is a key - yet just
unexplored - dimension to international mediation mechanisms in nuclear
disarmament processes. The observed empirical patterns across four
existing nuclear governances Mechanisms-Treaty on the Non-Proliferation
of Nuclear Weapons (NPT), International Atomic Energy Agency (IAEA)
safeguards regime, bilateral arms control like START/CTR and Central
Asian Nuclear-Weapon-Free Zone (CANWFZ)-demonstrate significant
differences in cost-effectiveness depending on particular elements of
institutional design. These implications directly challenge the
perception that more complex multilateral institutions generate better
outcomes and highlight how simpler, strategically designed agreements
often deliver as good or better outcomes at much lower cost.

The dissenting four main research hypotheses have been tested by the
documentary analysis and cost-structure comparison. First, hybrid
mediation models, whereby formal institutional mechanisms are integrated
with specialized technical cooperation instruments yield very high
cost-effectiveness ratios: more than 20:1 when measured in security
terms (see the START/CTR case). In this paper, hybridization
specifically denotes «nesting» technical cooperation and verification
toolkits inside treaty-based governance cycles (rather than running
parallel structures), which is the key mechanism behind the observed
cost-effectiveness differences. Second, efficient verification -
transparent systems the second point is that transparency in
verification can lead to much greater efficiency of mediation integrated
verification than parallel, autonomous systems of verification as shown
by IAEA IS analysis. Third, regional institutions such as the CANWFZ
also offer much greater cost-effectiveness than universal mechanisms
(through economies of scope, a minimum of overhead and management burden
sharing across states) as well as leveraged use of existing verification
infrastructure. Last one, institutional design attributes that focus on
data standardization and inspection protocol sharing effectively reduces
transaction costs of compliance verification.

This study contributes in important ways to our knowledge of
international governance and arms control. At a concept level, it
extends institutionalist concepts to incorporate transaction cost
economics constructs at the conceptual level, dumping empirics that
illustration was available of the extent design variables make for
differential cost-efficiency outcomes irrespective of de jure legal
authority. Disciplinary framing at the disciplinary level, the article
responds to a long-standing analytical divide between international law
scholarship (focused on verification standards and legal compliance) and
economics literature (analyzing implementation costs) by integrating
these insights in an overarching analytical framework, the cataloging of
individual design elements - host-country burden-sharing, integration
verification, standardization of data, and regional coordination -
yields operational advice for architects who build institutions in
future disarmament agreements.

Kazakhstan' s institutions of innovation will continue to
offer us evidence-based replicable models for cost-efficient disarmament
diplomacy, including voluntary denuclearization, CANWFZ leadership, IAEA
LEU Bank hosting and the International Day Against Nuclear Testing. The
documentary record shows that Kazakhstan attained these results not
through outsized financial resources but by strategic institutional
placement, moral authority based on historical fact and an efficient
burden-sharing arrangement. This is especially relevant for developing
countries that aim to shape international nuclear governance, but must
work with limited resources.

From a humanitarian perspective, the research uncovers an essential
economic irrationality in today's global security: countries spend
jointly some 750 times the annual verification budget of IAEA on
military expenditure with verifying infrastructure remaining chronically
underfunded - an estimate that carries distinct relevance to future
challenges posed by unilateral interests and approaches. This
misdirection is both a security problem and an economic opportunity.
Lifting even relatively small percentages of military expenditure from
weapons production for verification and mediation measures would vastly
increase disarmament enforcement capabilities. The remarkably favorable
cost benefit ratios found in this study -especially the 1:20 ratio of
return on investment for START/CTR agreements - indicate that
reallocating such funds would improve not only our security, but also
generate economic benefits.

This study suggests some potential areas for further study. The first is
that cost-effectiveness longitudinally analyzed as the Treaty on the
Prohibition of Nuclear Weapons (TPNW) takes hold and evidence is
collected could be useful for comparison purposes. Second, quantifying
the effects of innovations in verification technologies (for example:
zero-knowledge protocols and machine learning) might be able to
extrapolate cost savings expected from future disarmament verification
regimes. Third, cross-hybrid learning on burden-sharing mechanisms with
other international regimes (climate and humanitarian intervention and
pandemic governance) could potentially test the generalizability of
efficiency principles found in nuclear domains. Fourth, the cooperation
of institutional economics, technology assessment and analysis of
diplomatic negotiations could be utilized in interdisciplinary research
aimed at creating predictive models for the optimal solutions to
disarmament arrangement.

In summary, this study defines economic efficiency as an instrumental
and empirical facet of the institutional effectiveness in international
nuclear governance. By showing stem-levels that certain institutional
designs are more likely to generate cost-effective outcomes, the article
lays empirical groundwork for evidence-based institution design in
future disarmament negotiations. Developing countries have much to learn
from Kazakhstan' s demonstrated ability to drive globally
significant denuclearization initiatives and leverage resource
efficiencies, one that can be tailored for their own use. Its larger
relevance lies not only as a contribution to our knowledge, as it is an
outstanding model of how «real» disarmament diplomacy could work if
structured well or implemented wisely, but also it sheds light on the
fact that economic efficiency and disarmament effectiveness are not
mutually exclusive goals; rather they can be made compatible goals
through good institutional design. Such efficiency principles might well
be an interesting lens through which to focus on the institutional
design process in contributing to a greater likelihood of success with
disarmament agreement and more appropriate use, in conformity with the
logic of collective action, for scarce global resources committed to
nuclear governance/peace building priorities.

\emph{{\bfseries Funding.} This research is funded by the Science Committee
of the Ministry of Science and Higher Education of the Republic of
Kazakhstan, grant number AP 25795633.}

{\bfseries References:}

1. Treaty on the Non-Proliferation of Nuclear Weapons//International
Atomic Energy Agency. Vienna: IAEA Publications 1970. {[}E-resource{]}.
URL:\ul{\url{https://www.iaea.org/topics/non-proliferation-treaty}.}
Date of access: 27.11.2025.

2. Dorn W. The Case for a United Nations Verification Agency //
Disarmament Forum. - 2006 {[}E-resource{]}. URL:
https://walterdorn.net/20-the-case-for-a-united-nations-verification-agency-disarmament-under-effective-international-control.
Date of access: 27.11.2025

3. The NPT and the TPNW: Compatible or conflicting nuclear disarmament
agendas? // Stockholm International Peace Research Institute (SIPRI). -
2019. {[}E-resource{]}. URL:
https://www.sipri.org/commentary/blog/2019/npt-and-tpnw-compatible-or-conflicting-nuclear-weapons-treaties.
Date of access: 27.11.2025

4. Speakers Urge Efficient Financial Management, Avoiding Duplicate
Processes, as Fifth Committee Examines Programme Budget Implications of
Disarmament-Related Texts // UN Press Release. - 2023. {[}E-resource{]}.
URL:
\href{https://press.un.org/en/2023/gaab4445.doc.htm}{\ul{https://press.un.org/en/2023/gaab4445.doc.htm}}
Date of access: 27.11.2025

5. Susan W. Costs of Disarmament --- Disarming the Costs: Nuclear Arms
Control and Nuclear Rearmament // UNIDIR. - 2003. {[}E-resource{]}. URL:
\href{https://unidir.org/files/publication/pdfs/costs-of-disarmament-disarming-the-costs-nuclear-arms-control-and-nuclear-rearmament-306.pdf}{\ul{https://unidir.org/files/publication/pdfs/costs-of-disarmament-disarming-the-costs-nuclear-arms-control-and-nuclear-rearmament-306.pdf}}
Date of access: 27.11.2025

6. Central Asian Nuclear-Weapon-Free-Zone // NTI. {[}E-resource{]}. URL:
\href{https://www.nti.org/education-center/treaties-and-regimes/central-asia-nuclear-weapon-free-zone-canwz/}{\ul{https://www.nti.org/education-center/treaties-and-regimes/central-asia-nuclear-weapon-free-zone-canwz/}}
Date of access: 27.11.2025

7. Abuova N. Kazakhstan Hosts Workshop to Strengthen Cooperation between
Nuclear-Free Zones// Astana Times. - 2024. {[}E-resource{]}. URL:
https://astanatimes.com/2024/08/kazakhstan-hosts-workshop-to-strengthen-cooperation-between-nuclear-free-zones/
Date of access: 27.11.2025

8. Kazakhstan's contribution to nuclear disarmament, non-proliferation
and peaceful use of nuclear energy // Government of Kazakhstan (Ministry
of Foreign Affairs). - 2021. {[}E-resource{]}. URL:
https://www.gov.kz/memleket/entities/mfa-amman/press/article/details/59812?lang=en.-
Date of access: 04.01.2026.

9. Kreps S.E. The Institutional Design of Arms Control Agreements //
Foreign Policy Analysis. - 2018. - Vol.14(1) - P.127-147. DOI
\href{https://doi.org/10.1093/fpa/orw045}{10.1093/fpa/orw045}.

10. Fortna V. P. Peace Time: Cease-Fire Agreements and the Durability of
Peace // Princeton University Press. - 2004.- 264 p.

11. Baumol W. J. Williamson's The Economic Institutions of Capitalism
{[}Review of The Economic Institutions of Capitalism., by O. E.
Williamson{]} // The RAND Journal of Economics. - 1986. - 17(2). - P.
279-286. DOI 10.2307/2555390.

12. Fuhrmann M. and Yonatan Lupu. Do Arms Control Treaties Work?
Assessing the Effectiveness of the Nuclear Nonproliferation Treaty //
International Studies Quarterly. - 2016. - 60, - No.3. - P.530-539.
DOI \href{https://doi.org/10.1093/isq/sqw013}{10.1093/isq/sqw013}.

13. Erästö Tytti, Ugnė Komžaitė and Petr Topychkanov. Operationalizing
nuclear disarmament verification// Stockholm international peace
research institute. - 2019. {[}E-resource{]}. URL:

https://www.sipri.org/sites/default/files/2019-04/sipriinsight1904\_0.pdf.-
Date of access: 27.11.2025.

14. Zhomartkyzy M. The role of mediation in international conflict
resolution // Law and Safety. - 2023. - 90(3). - P.169-178.
\href{https://doi.org/10.32631/pb.2023.3.14}{DOI 10.32631/pb.2023.3.14}.

15. The Agency' s Draft Budget Update for 2025. Document
GC/68-5 // Vienna: IAEA. {[}E-resource{]}. URL:
\ul{\url{https://www.iaea.org/sites/default/files/gc/gc68-5.pdf}. -}
Date of access: 27.11.2025.

16. Mayhew N., Kirsten I. Navigating the IAEA Budget Process // Vienna
Center for Disarmament and Non-Proliferation. - 2023. {[}E-resource{]}.
URL:
\href{https://vcdnp.org/wp-content/uploads/2023/08/print_navigating_the_iaea_budget_process_rev.1.pdf}{\ul{https://vcdnp.org/wp-content/uploads/2023/08/print\_navigating\_the\_iaea\_budget\_process\_rev.1.pdf}}
.-Date of access: 27.11.2025.

17. SIPRI contributes to global UN report on military expenditure //
Stockholm: Stockholm International Peace Research Institute. - 2025.
{[}E-resource{]}. URL:
\ul{\url{https://www.sipri.org/news/2025/sipri-contributes-global-un-report-military-expenditure}.-}
Date of access: 27.11.2025

18. Ileana Exaras. Military spending worldwide hits record \$2.7
trillion // UN News. - 2025. {[}E-resource{]}. URL:
\ul{\url{https://news.un.org/en/story/2025/09/1165809}.} - Date of
access: 27.11.2025.

19. Record military spending threatens global peace and development, new
UN report warns// New York: United Nations Development Programme. -
2025. {[}E-resource{]}. URL:
\ul{https://www.undp.org/press-releases/record-military-spending-threatens-global-peace-and-development-new-un-report-warns.-}
Date of access: 27.11.2025.

20. START I. // US Department of state. - 2004 {[}E-resource{]}. URL:
\href{https://www.state.gov/2004-TIAS/}{2004 Treaties and Agreements -
United States Department of State} .-Date of access: 27.11.2025

21. From Donor to Partner: The Evolution of U.S. Cooperative Threat
Reduction // Washington: Nuclear Threat Initiative. - 2011.
{[}E-resource{]}. URL:
\ul{\url{https://www.nti.org/analysis/articles/donor-partner-evolution-us-cooperative-threat-reduction-global-security-engagement/}.-}
Date of access: 27.11.2025

22. Cooperative Threat Reduction Program Fiscal Year (FY) 2025 Budget
Estimates // Washington: DoD Comptroller. - 2024 {[}E-resource{]}. URL:
https://comptroller.war.gov/Portals/45/Documents/defbudget/FY2025/budget\_justification/pdfs/01\_Operation\_and\_Maintenance/O\_M\_VOL\_1\_PART\_2/CTR\_OP-5.pdf.-
Date of access: 27.11.2025.

23. Kakatkar Manasi and Miles A. Pomper. Central Asian
Nuclear-Weapon-Free Zone Formed// Arms Control Today. - 39. - No.1. -
2009. - P.42 - 42.

24. Marco Roscini. Something Old, Something New: The 2006 Semipalatinsk
Treaty on a Nuclear Weapon-Free Zone in Central Asia // Chinese Journal
of International Law. -Vol.7(3). - 2008. - P.593-624. DOI
\href{https://doi.org/10.1093/chinesejil/jmn033}{10.1093/chinesejil/jmn033}.

25. NPT Review Process: An Explainer // Vienna: Vienna Center for
Disarmament and Non-Proliferation. - 2023 {[}E-resource{]}. URL:
\url{https://vcdnp.org/wp-content/uploads/2023/07/web_npt_review_process_07212023.pdf}
. -Date of access: 27.11.2025.

26. Financing the NPT Review Process. Working paper presented to the
second session of the Preparatory Committee for the 2010 Review
Conference // US Department of State. - 2008 {[}E-resource{]}. URL:
\url{https://2001-2009.state.gov/t/isn/rls/other/114114.htm}. - Date of
access: 27.11.2025.

27. Zoe N. Gastelum. Quantifying Safeguards Innovations: Impacts of
Innovations on International Atomic Energy Agency Safeguards Inspection
Effort // Journal for Peace and Nuclear Disarmament. - 2024. - No 7(1).
- P.236-263. - DOI 10.1080/25751654.2024.2364962.

28. Verification of Nuclear Disarmament: Capstone Report // Washington:
International Partnership for Nuclear Disarmament Verification. - 2024
{[}E-resource{]}. URL:
\ul{\url{https://www.ipndv.org/wp-content/uploads/2024/06/IPNDV-Capstone_FINAL-1.pdf}.-}
Date of access: 27.11.2025.

29. Philippe S., Goldston R., Glaser A. et al. A physical zero-knowledge
object-comparison system for nuclear warhead verification // Nat Commun.
-2016.- No 7:12890. \href{https://doi.org/10.1038/ncomms12890}{DOI
10.1038/ncomms12890}.

30. D. Kairolda. The Silent Fallout: The Legacy of Soviet Nuclear
Testing in Kazakhstan // Nuclear Age Peace Foundation. - 2025.
{[}E-resource{]}. URL:
\ul{\url{https://www.wagingpeace.org/the-silent-fallout-the-legacy-of-soviet-nuclear-testing-in-kazakhstan-and-its-modern-day-impacts/}-}
Date of access: 27.11.2025.

31. Satubaldina A. UN' s Izumi Nakamitsu Highlights
Kazakhstan' s Moral Leadership in Nuclear Disarmament //
The Astana Times. - 2024 {[}E-resource{]}. URL:
\ul{\url{https://astanatimes.com/2024/08/uns-izumi-nakamitsu-highlights-kazakhstans-moral-leadership-in-nuclear-disarmament/}.-}
Date of access: 27.11.2025.

32. Nuclear Disarmament Kazakhstan // Washington: NTI. - 2024.
{[}E-resource{]}. URL:
\url{https://www.nti.org/analysis/articles/kazakhstan-nuclear-disarmament/.-}
Date of access: 27.11.2025.

33. Velichkov K. Kazakhstan's nuclear governance as a foreign policy
ASSET // NNC RK Bulletin. - 2021. - No (3). -P.21-28.
\href{https://doi.org/10.52676/1729-7885-2021-3-21-28}{DOI
10.52676/1729-7885-2021-3-21-28}.

34. Treaty on the Prohibition of Nuclear Weapons // New York: United
Nations. - 2017 {[}E-resource{]}. URL:
\href{https://ihl-databases.icrc.org/assets/treaties/640-TPNW-EN.pdf}{\ul{https://ihl-databases.icrc.org/assets/treaties/640-TPNW-EN.pdf}}
.- Date of access: 27.11.2025.

35. Patton T., Philippe S., \& Mian Z. Fit for Purpose: An Evolutionary
Strategy for the Imple-mentation and Verification of the Treaty on the
Prohibition of Nuclear Weapons//Journal for Peace and Nuclear
Disarmament. - 2019. - No 2(2). - P.387 - 409.
\href{https://doi.org/10.1080/25751654.2019.1666699}{DOI
10.1080/25751654.2019.1666699}.

36. Amanzholov Zh. M., Ahmetov E. B. Bez"yadernaya zona v
Tsentral' noi Azii i ee rol'{} v ramkakh
voenno-politicheskogo izmereniya deyate// International Relations and
International 38. Law Journal. - 2016. - No.59(3). - S.42-48. {[}in
Russian{]}

37. Budlong J. The Institutional Design of Arms Control: To What Extent
Does Institutional Design Increase the Longevity of Arms Control
Agreements? // Electronic Theses and Dissertations of University of
Denver. - 2021. {[}E-resource{]}. URL:
\ul{\url{https://digitalcommons.du.edu/etd/1921}.-} Date of access:
27.11.2025.38. What are the most effective conflict mediation techniques used in
international diplomacy, and how can case studies from the United
Nations support these methods? // PsicoSmart. - 2025. {[}E-resource{]}.
URL:
\url{https://blogs.psico-smart.com/blog-what-are-the-most-effective-conflict-mediation-techniques-used-in-inte-188010}.
- Date of access: 27.11.2025.

39. Global spending on nuclear weapons topped \$100 billion in 2024 //
Geneva: International Campaign to Abolish Nuclear Weapons. - 2025.
{[}E-resource{]}. URL:
\ul{\url{https://www.icanw.org/global_spending_on_nuclear_weapons_topped_100_billion_in_2024}.-}
Date of access: 27.11.2025.

40. The Security We Need: Rebalancing Military Spending for a
Sustainable and Peaceful Future // United Nations: Peace and Security. -
2025. {[}E-resource{]}. URL:
\ul{\url{https://www.un.org/en/peace-and-security/the-true-cost-of-peace}.-}
Date of access: 27.11.2025.

\emph{{\bfseries Information about authors}}

Bolat A. -master of law, I. Zhansugurov Zhetysu University Taldykorgan,
Taldykorgan, Kazakhstan,

e-mail:
a.bolat.edu@gmail.com;

Kassymova A.- PhD, I. Zhansugurov Zhetysu University, Kazakhstan,
e-mail:
aiko\_1990kz@mail.ru;

Jangeldina D. - Candidate of Pedagogical Sciences, Associate Professor,
K. Kulazhanov Kazakh University of Technology and Business, Astana,
Kazakhstan, e-mail:
dariga\_da@mail.ru;

Kurmankulova N.- candidate of economical sciences, associate professor,
K. Kulazhanov Kazakh University of Technology and Busines, Astana,
Kazakhstan, e-mail:
nurgik-7@mail.ru.

\emph{{\bfseries Сведения об авторах}}

Болат А.А.- магистр юридических наук, Жетысуский университет им. И.
Жансугурова, Талдыкорган, Республика Казахстан, e-mail:
a.bolat.edu@gmail.com;

Касымова А.М.-доктор философии (PhD), Жетысуский университет им. И.
Жансугурова, Талдыкорган, Республика Казахстан, e-mail:
aiko\_1990kz@mail.ru;

Джангельдина Д.И.-к.п.н., доцент, Казахский университет технологии и
бизнеса им. К. Кулажанова, Астана, Қазақстан, e-mail:
dariga\_da@mail.ru;

Курманкулова Н.Ж.- к.э.н., асс.профессор Казахский университет
технологии и бизнеса им. К. Кулажанова, Астана, Қазақстан, e-mail:
nurgik-7@mail.ru.

МНРТИ 06.81.30

{\bfseries MERGERS AND ACQUISITIONS: THE ROLE OF DUE DILIGENCE IN M\&A
OUTCOMES}

{\bfseries A.
Amanbekova}\fig{e/image1}[]

\emph{Kazakh- British Technical University, Almaty, Kazakhstan}

\envelope Corresponding author:
amanbekovaidana@gmail.com

The article establishes due diligence as the essential factor which
decides M\&A transaction success through its core purpose. Organizations
need due diligence (hereinafter referred to as ``DD'') as their
essential M\&A success factor because it enables them to reach their
objectives through unstable market conditions and complicated business
operations. The Bain \& Company survey shows that most failed deals
result from various factors which include DD-related issues that lead to
incorrect synergy assessments and failed due diligence that missed
important problems and target companies that received artificial
enhancements for sale purposes and targets that lack proper strategic
alignment and teams ignored their concerns while participating in the
competitive auction process. The research findings show that DD serves
as the fundamental component because it establishes the level at which
organizations achieve their strategic decision implementation regarding
integration and synergy assessment and risk management.

The research uses secondary data analysis to develop multiple proxy
variables which measure DD quality and then evaluates these variables
against their effects on deal achievement. The regression analysis is
performed to assess outcomes by using the change in acquirer stock price
as the dependent variable which measures price changes from
pre-announcement to post-closing periods.

{\bfseries Keywords:} due diligence, mergers and acquisitions, role of due
diligence in deal success, stock price change after deals, market
volatility, regression analysis, share price reaction

{\bfseries БІРІГУ ЖӘНЕ ЖҰТЫЛУ: БІРІГУ ЖӘНЕ ЖҰТЫЛУ НӘТИЖЕЛЕРІНДЕГІ ДЬЮ
ДИЛИДЖЕНСТІҢ РӨЛІ}

{\bfseries А. М. Аманбекова\envelope }

\emph{Қазақстан-Британ Техникалық Университеті, Алматы, Қазақстан,}

e-mail:
amanbekovaidana@gmail.com

Бұл мақаланың мақсаты - бірігу және жұтылу мәмілелеріндегі дью дилидженс
процесінің маңызын түсіндіру. Нарықтағы жоғары құбылмалылық және
корпоративтік стратегиялардың барған сайын күрделенуі жағдайында дью
дилидженс (ДД) бірігу және жұтылу (M\&A) мәмілелерінің тиімділігін
қамтамасыз етудегі негізгі кезеңге айналады. Bain \& Company жүргізген
сауалнама нәтижелеріне сәйкес, сәтсіз мәмілелердің көпшілігі бірқатар
себептермен түсіндіріледі, олардың көбі ДД-ге қатысты факторлар болып
табылады: синергияның асыра бағалануы, дью дилидженстің негізгі
мәселелерді анықтай алмауы, мақсатты компанияның сату үшін «әсемделуі»,
стратегиялық сәйкестіктің жеткіліксіздігі, күмәнді белгілердің еленбеуі,
сондай-ақ тендерлік бәсекенің қызуына еріп кету. Осы тұжырымдарға сүйене
отырып, ДД рөлін іргелі деп санауға болады: дәл осы кезеңде интеграция,
синергияны бағалау және тәуекелдерді басқару бойынша кейінгі
стратегиялық шешімдердің сапасы қалыптасады.

Бұл жұмыста екінші реттік деректерді талдау тәсілі қолданылады: ДД
сапасын өлшеу үшін бірнеше прокси-көрсеткіш ұсынылып, олардың мәміленің
табыстылығына әсерін эмпирикалық түрде тексеріледі. Нәтижелерді бағалау
үшін регрессиялық талдау жүргізіледі, мұнда тәуелді айнымалы - мәміле
жарияланғанға дейінгі кезең мен мәміле жабылғаннан кейінгі кезең
арасындағы сатып алушы компанияның акция бағасының өзгерісі.

{\bfseries Түйін сөздер:} дью дилидженс, бірігу және жұтылу, мәміленің
сәтті болуындағы дью дилидженстің рөлі, мәмілелерден кейін акция
бағаларының өзгеруі, нарық құбылмалылығы, регрессиялық талдау, акция
бағаларының реакциясы

\begin{quote}
{\bfseries СЛИЯНИЯ И ПОГЛОЩЕНИЯ: РОЛЬ ДЬЮ ДИЛИДЖЕНСА В РЕЗУЛЬТАТАХ СЛИЯНИЙ
И ПОГЛОЩЕНИЙ}
\end{quote}

{\bfseries А. М. Аманбекова\envelope }

\emph{Казахстанско-Британский Технический Университет, Алматы,
Казахстан,}

e-mail:
amanbekovaidana@gmail.com

Цель этой статьи - объяснить важность дью дилидженса в сделках слияний и
поглощений. В условиях высокой волатильности рынков и усложнения
корпоративных стратегий дью дилидженс (ДД) становится ключевым этапом
обеспечения эффективности сделок по слияниям и поглощениям (M\&A).
Согласно результатам опроса Bain \& Company, большинство неудачных
сделок объясняется рядом причин, многие из которых являются
ДД-факторами, такими как: переоценка синергий, неспособность дью
дилидженса выявить ключевые проблемы, «приукрашивание» целевой компании
перед продажей, недостаточное стратегическое соответствие, игнорирование
сомнений и предупреждающих сигналов, а также вовлечение в азарт
конкурентных торгов. Исходя из этих выводов, роль дью дилидженса можно
считать фундаментальной: именно на этом этапе формируется качество
дальнейших стратегических решений по интеграции, оценке синергий и
управлению рисками.

В данной статье используется подход анализа вторичных данных,
предлагающий ряд прокси-показателей для измерения качества ДД и
эмпирически оценивающий их влияние на успешность сделки. Для оценки
результатов проводится регрессионный анализ, где зависимой переменной
выступает изменение стоимости акций компании-покупателя между периодом
до объявления сделки и периодом после её закрытия.

{\bfseries Ключевые слова:} дью дилидженс, слияния и поглощения, роль дью
дилидженса в успешности сделок, изменение цены акций после сделок,
волатильность рынка, регрессионный анализ, реакция цены акций

{\bfseries Introduction.} Mergers and acquisitions (hereinafter referred to
as ``M\&A'') require due diligence (hereinafter referred to as ``DD'')
because M\&A transactions generate unpredictable results which result in
various acquisition deal outcomes. M\&A serve as a strategic growth tool
which appears frequently, and many research show that acquisition
performance for acquiring companies remains uncertain after deal
completion.

For example, a key study by King et al. (2003) {[}1{]} summarizes many
empirical works, which prove that in the average M\&A deals results are
not consistently positive for acquirers, and the variation in outcomes
remains large and partly unexplained. Researchers continue to seek
``moderator factors'' which will help them predict deal success or
failure during the initial preparation phase. Due diligence is important
here, because it is a systematic review of the target company that helps
identify risks, assumptions, and structural issues that can strongly
affect post-deal performance.

Market practice confirms that due diligence is one of the main
``bottlenecks'' in successful deal closing. The KPMG 2025 M\&A Deal
Market Study (survey of 300 deal professionals) {[}2{]} shows that
completing due diligence is among the top three obstacles to closing
deals, in line with valuation and integration-related issues. The study
reports that 41\% of respondents mention completing due diligence as a
critical difficulty during the deal process. This supports the idea that
due diligence is not a formal checklist, but a key stage that can slow
down, complicate, or block a deal. Also, according to this study, the
rising number of deal risks demands advanced valuation models which need
more information and stronger regulatory standards, thus due diligence
serves as a vital assessment tool that affects deal length, market
confidence, and transaction success rates.

Consulting practice also indicates that most M\&A failures result from
integration problems and incorrect synergy assessments, which should be
detected during due diligence evaluation. For example, the Bain report
``The Merger Before the Merger'' {[}3{]} lists key reasons why deals
break down, according to their conducted survey, which include (Fig.1):

\begin{quote}
- Due diligence failed to highlight key issues

- Target was dressed up for sale

- Insufficient strategic fit

- Doubts were pushed aside

- Caught up in the bidding process
\end{quote}

These factors are strongly linked to due diligence quality. They create
``hidden risk'' effects where the buyer overestimates value,
underestimates integration challenges, and closes the deal with
incomplete information.

\fig{e/image9}[]

{\bfseries Fig.1 - Results of the Bain \& Company Survey (n = 250)}

In addition, the McKinsey report "Where mergers go wrong" {[}4{]} also
identifies overestimation of synergies as a stable problem which results
in the "winner' s curse" because the auction winner ends
up paying excessive amounts as a result of unrealistic expectations and
insufficient proof of their assumptions. The process fails to
demonstrate actual project risks and integration challenges and required
resources when due diligence is not fully completed.

Hence, the research findings will have the following practical
significance:

- due diligence factors linked to deal success can transform due
diligence from ``checklist'' to measurable KPIs and predictive risk
markers. If due diligence quality indicators statistically explain
buyer's stock price reaction to the conducted deals, then they can be
translated into management actions as: 1) Due diligence scope and depth
requirements;

- criteria for advisor selection;

- risk adjusted valuation and SPA protection.

The theoretical significance of the research will combine three
literature streams that usually studied separately:

- event studies (market reaction to corporate events);

- accounting signals (post-deal reporting outcomes, such as impairment);

- due diligence literature (due diligence as a tool to reduce
information asymmetry).

Finally, successful M\&A is important not only on the corporate level,
but beyond it. M\&A reallocates capital, restructures industries, and
transforms economies. Research suggests deals shape market structure,
enable exit of weak players, speed up technology renewal, and help
industries adapt to shocks (Andrade, Mitchell \& Stafford, 2001 {[}5{]};
Andrade, 2004 {[}6{]}; Breinlich, 2008 {[}7{]}). Cross-border M\&A is
also part of FDI flows and supports transfer of technology and
management practices (UNCTAD WIR series {[}8{]}). Failures destroy value
and may create political risks and lower trust in cross-border
investments.

In addition, the effectiveness of a due diligence can vary depending on
the mechanism for the purchase price under the Sales and Purchase
Agreement (SPA). In international M\&A practice, the two most common
mechanisms are the ``locked box'' and ``completion accounts''. Due
diligence under the two mechanisms is carried out differently.

Under the locked box mechanism, the purchase price is determined based
on the company's historical financial statements as of a particular date
in the past (the ``locked box date''). This means that the due diligence
exercise will have two phases: 1) a detailed review of the historical
financial statements and the key financial metrics underlying the
purchase price calculation, and 2) a more general review of the period
since the locked box date to ensure that there have been no material
leakage transactions which could allow the sellers to extract value from
the company at the buyer's expense.

On the other hand, under the completion accounts mechanism, the purchase
price is calculated on the basis of the financial performance of target
company as of the date of completion of transaction. Therefore, the
price adjustment is done post-closing and the need of detailed review of
the period between signing and closing of the transaction is reduced.
The target company's historical financial statements and financial
metrics, such as net debt and net working capital, are reviewed and are
further included in the purchase price calculation for the adjustment
for completion accounts.

Thus, the due diligence under locked box mechanism is generally more
detailed, since the purchase price is fixed, and a historical account
has to be verified as to its accuracy and adequacy for the period from
the effective date to the closing price. In a completion accounts
mechanism, the focus is on the historical aspects and the contingent
items which will impact future price adjustments, such as net debt and
net working capital.

{\bfseries Materials and мethods.} This research examines whether higher
due diligence quality leads to improved post-deal performance for the
acquiring company. The central hypothesis proposes that a higher level
of due diligence quality has a positive impact on the acquirer's
subsequent performance.

Research design and stages. The research proceeds through several steps.
The first step was collecting secondary deal and market data, after that
I built DD quality proxies and other control variables. After all data
was collected, I estimated a linear regression model and checked
multicollinearity using VIF. Lastly, I interpreted coefficients and
statistical significance.

Data and materials. For the purposes of this research, I used secondary
data collected from CapitalIQ and Refinitiv. The sample includes 39
deals, and sample size adequacy is justified using MacKinlay (1997)
{[}9{]} power analysis. In MacKinlay's example, event study test power
with sample size around 30 is about 0.54, which has the power to
identify a real effect. Since this study uses 39 observations (larger
than 30), expected statistical power should not be lower than in that
illustration. The research design with this sample size allows
regression analysis through event study methods which follow established
requirements.

Variables. The dependent variable measures post-deal performance for the
acquirer and taken as a stock price change between before deal
announcement and one year after deal closing (long-run returns). This
reflects realization of expectations, ability to achieve synergies,
integration success, and avoidance of overpayment. This approach is
supported by long-run return studies: 1) Loughran \& Vijh (2012)
{[}10{]}: long-run returns of acquirers differ from short-run returns
and depend on deal form and integration dynamics; 2) Rau \& Vermaelen
(1998) {[}11{]}: post-deal performance differs by bidder type and deal
characteristics.

The \emph{independent variables} are DD quality proxies (Table 1). A key
limitation is that direct DD scope and true DD quality are not
observable, as there is no available information about its underlying
scope. Therefore, I used proxy variables:

- \emph{information complexity of the target} - target type (private /
public), industry (whether target is in a different industry than the
buyer), country (cross-border or not);

- \emph{deal duration} - days between announcement and closing. This is
not a direct DD metric, but assumed that more time can allow more
complete DD (based on {[}12{]});

- \emph{deal advisor} - advisor quality is treated as a valid
independent variable because top-tier M\&A advisors can provide better
analysis, better risk detection, and better structuring, affecting both
short-run market reaction and long-run performance. This is supported by
work linking advisor choice to bidder returns (Golubov, Petmezas \&
Travlos {[}13{]}; Rau, 2000 {[}14{]}). The mechanism is reduction of
information asymmetrical and lower risk of overpayment or wrong synergy
assumptions;

- \emph{goodwill impairment} - goodwill impairment in the first 1-3
years after the deal is used as a proxy for DD quality because it
reflects how correct pre-acquisition assessments were, whether expected
synergies were realized, and whether acquirer overpaid for the target
company. The research by Wangerin (2019) {[}12{]} and Potepa \& Thomas
(2022) {[}15{]} shows that poor debt disclosure practices result in
higher impairment costs and lower business performance and wrong fair
value evaluations. Therefore, impairment is treated as an objective
ex-post signal of DD effectiveness.

Lastly, controls include payment method, acquirer size and performance,
and relative deal size.

%% \begin{longtable}[]{@{}
%%   >{\raggedright\arraybackslash}p{(\linewidth - 6\tabcolsep) * \real{0.2319}}
%%   >{\centering\arraybackslash}p{(\linewidth - 6\tabcolsep) * \real{0.1419}}
%%   >{\raggedright\arraybackslash}p{(\linewidth - 6\tabcolsep) * \real{0.4171}}
%%   >{\raggedright\arraybackslash}p{(\linewidth - 6\tabcolsep) * \real{0.2090}}@{}}
%% \toprule\noalign{}
%% \begin{minipage}[b]{\linewidth}\centering
%% \emph{{\bfseries Variable name}}
%% \end{minipage} & \begin{minipage}[b]{\linewidth}\centering
%% \emph{{\bfseries Unit}}
%% \end{minipage} & \begin{minipage}[b]{\linewidth}\centering
%% \emph{{\bfseries Definition}}
%% \end{minipage} & \begin{minipage}[b]{\linewidth}\centering
%% \emph{{\bfseries Source}}
%% \end{minipage} \\
%% \midrule\noalign{}
%% \endhead
%% \bottomrule\noalign{}
%% \endlastfoot
%% Performance & ln & Natural logarithm of the acquirer' s
%% stock price change (pre-announcement vs 1 year after closing) &
%% CapitalIQ, Refinitiv \\
%% Target type & dummy & Whether the target company is private or public &
%% CapitalIQ, Refinitiv \\
%% Industry & dummy & Whether the target company from the same industry as
%% acquirer & CapitalIQ, Refinitiv \\
%% Country & dummy & Is the deal cross-border & CapitalIQ, Refinitiv \\
%% Deal duration & days & Number of days between deal announcement and deal
%% closing & CapitalIQ, Refinitiv \\
%% Advisor & dummy & Whether the deal advisor was top-tier (e.g. Morgan
%% Stanley, Goldman Sachs etc.) & CapitalIQ, Refinitiv \\
%% Impairments & dummy & Wheter the acquirer had a goodwill impairment in
%% the first 1-3 years after the deal closed & CapitalIQ, Refinitiv \\
%% Payment & dummy & The deal' s payment method (cash /
%% equity) & CapitalIQ, Refinitiv \\
%% Relative deal size & USDm & The deal' s size relatively
%% to acquirer' s market capitalization & CapitalIQ,
%% Refinitiv \\
%% Acquirer equity value & USDm & The equity value of acquirer & CapitalIQ,
%% Refinitiv \\
%% Acquirer ROA & \% & ROA of acquirer & CapitalIQ, Refinitiv \\
%% Acquirer leverage & \% & Acquirer' s debt-to-equity ratio
%% & CapitalIQ, Refinitiv \\
%% \end{longtable}

\emph{{\bfseries Table 1 - Variables for regression}}

Regression model. For the purposes of this research I made multifactor
linear regression model:

Performanceᵢ = β₀ + β₁·DD Qualityᵢ + β₂·Controlsᵢ + εᵢ

\begin{quote}
Where: 1) Performanceᵢ = stock price change (pre-announcement vs 1 year
after closing); 2) DDQualityᵢ = DD proxies; 3) Controlsᵢ = control
variables; 4) εᵢ = error term
\end{quote}

After regression model was runed, I also conducted multicollinearity
test, where VIF = 1.432, which suggests no multicollinearity concerns.

{\bfseries Results and Discussion.} In this study, I evaluate statistical
significance at the 90\% confidence level (α = 0.10). This choice is
based on the event-study econometrics literature, which explains that
return-based tests can have limited statistical power in small samples
and that inference depends on sample size and return volatility.
Following this guidance, I report results that are significant at the
10\% level as economically meaningful signals that should be tested
further in larger samples. (Kothari \& Warner, 2004) {[}16{]}.

Model fit. The model explains about 30\% of the variation in the
dependent variable (R² = 0.302), while the adjusted R² is 0.140 (Table
2). In M\&A return regressions, such values are acceptable because stock
return outcomes are noisy and influenced by many unobservable factors.
Prior evidence suggests that cross-sectional models predicting acquirer
stock-return reactions often have low explanatory power, so these fit
measures are reasonable for this market setting. (Quariguasi Frota Neto
et al., 2025 {[}17{]}).

{\bfseries Table 2 - Multifactor regression model results}

%% \begin{longtable}[]{@{}
%%   >{\raggedright\arraybackslash}p{(\linewidth - 2\tabcolsep) * \real{0.5676}}
%%   >{\raggedleft\arraybackslash}p{(\linewidth - 2\tabcolsep) * \real{0.4324}}@{}}
%% \toprule\noalign{}
%% \multicolumn{2}{@{}>{\centering\arraybackslash}p{(\linewidth - 2\tabcolsep) * \real{1.0000} + 2\tabcolsep}@{}}{%
%% \begin{minipage}[b]{\linewidth}\centering
%% \emph{{\bfseries Regression Statistics}}
%% \end{minipage}} \\
%% \midrule\noalign{}
%% \endhead
%% \bottomrule\noalign{}
%% \endlastfoot
%% Multiple R & 0.549265181 \\
%% R Square & 0.301692239 \\
%% Adjusted R Square & 0.139509534 \\
%% Standard Error & 0.415057657 \\
%% Observations & 39 \\
%% \end{longtable}

Key DD-related findings (Table 3). Advisor (dummy) has a positive
coefficient (+0.344) and is marginally significant (p ≈ 0.07). The
coefficient on Advisor is +0.344 with p = 0.070. Since p <{}
0.10, this variable is significant at the 90\% confidence level. The
positive sign indicates that working with an experienced (top-tier)
advisor who represents superior deal due diligence and structure leads
to better one-year post-deal stock performance for the acquiring
company.

{\bfseries Table 3 - Multifactor regression model results}

%% \begin{longtable}[]{@{}
%%   >{\raggedright\arraybackslash}p{(\linewidth - 6\tabcolsep) * \real{0.3206}}
%%   >{\raggedleft\arraybackslash}p{(\linewidth - 6\tabcolsep) * \real{0.2166}}
%%   >{\raggedleft\arraybackslash}p{(\linewidth - 6\tabcolsep) * \real{0.2532}}
%%   >{\raggedleft\arraybackslash}p{(\linewidth - 6\tabcolsep) * \real{0.2096}}@{}}
%% \toprule\noalign{}
%% \begin{minipage}[b]{\linewidth}\centering
%% \emph{~}
%% \end{minipage} & \begin{minipage}[b]{\linewidth}\centering
%% \emph{{\bfseries Coefficients}}
%% \end{minipage} & \begin{minipage}[b]{\linewidth}\centering
%% \emph{{\bfseries Standard Error}}
%% \end{minipage} & \begin{minipage}[b]{\linewidth}\centering
%% \emph{{\bfseries P-value}}
%% \end{minipage} \\
%% \midrule\noalign{}
%% \endhead
%% \bottomrule\noalign{}
%% \endlastfoot
%% Target type (dummy) & 0.076671579 & 0.165979733 & 0.647254323 \\
%% Industry (dummy) & 0.031494975 & 0.14875856 & 0.833669167 \\
%% Country (dummy) & -0.097429955 & 0.186519835 & 0.605019448 \\
%% Deal duration & -3.94365E-05 & 0.001044216 & 0.970108425 \\
%% Advisor (dummy) & 0.344301484 & 0.183779515 & 0.070164574 \\
%% Impairments (dummy) & -0.301884692 & 0.154618128 & 0.059679329 \\
%% Payment (dummy) & 0.101490258 & 0.153903257 & 0.5143302 \\
%% \end{longtable}

Impairments (dummy) has a negative coefficient (-0.302) and is also
marginally significant (p ≈ 0.06) at the 90\% confidence level. The
negative sign indicates that deals followed by goodwill impairment
within 1-3 years tend to have worse one-year post-deal stock price
performance. The evidence supports impairment as an ex-post signal which
reveals unmet expectations and possible overpayment and unidentifiable
risks that became apparent after the transaction completion.

Other variables. The regression results show that target type and
industry similarity and country and deal duration and payment method do
not affect the outcome because their p-values are high. So in this
sample and this specification they do not show strong independent
explanatory power.

{\bfseries Conclusion.} The goal of this study is to test whether due
diligence quality affects M\&A outcomes for the acquirer. The study uses
a regression model within an event-study logic, where performance is
measured as the acquirer's stock price change from before the
announcement to one year after deal closing. DD quality is proxied by
deal complexity measures, deal duration, advisor dummy, and post-deal
goodwill impairment. The results demonstrate the following findings. The
analysis shows that the Advisor dummy variable demonstrates a positive
relationship with post-deal performance results at a marginally
significant level (p ≈ 0.07), the Goodwill impairment dummy variable
produces negative results which affect post-deal performance at a
marginally significant level (p ≈ 0.06). Other variables are not
significant. The present sample needs no further control procedures for
the current assessment, as VIF indicates no multicollinearity (VIF ≈
1.43).

The research findings indicate that DD-related conditions create
significance for what happens after a deal is completed. Specifically:
1) Better advisor-related conditions (proxy for stronger analysis and
structuring) are associated with higher post-deal stock price
performance; 2) Goodwill impairment in the first 1-3 years after a deal
is associated with worse post-deal performance and may reflect weak
pre-deal screening, weak DD, or unrealistic synergy/valuation
assumptions.

So, as a result, the research results confirm that due diligence
functions as an invisible factor which links expected deal outcomes to
actual post-acquisition performance.

{\bfseries References}

1. King D.R., Dan D.R., Daily C.M., Covin J.G. Meta-analyses of
post-acquisition performance: indications of unidentified moderators//
Strategic Management Journal. -2004. - Vol.25(2). - P.187-200. DOI
10.1002/smj.371.

2. The KPMG 2025 M\&A Deal Market Study. - 2025. URL:
\url{https://kpmg.com/us/en/articles/2025/2025-ma-deal-market-study.html.-}
Date of access: 03.01.2026

3. Grob J., Meacham M. The merger before the merger/Bain and Company.
-1999. . URL:
https://www.bain.com/contentassets/afa083a686c847d58c7fef8e6d0ef379/bb\_merger\_before\_the\_merger.pdf
.-Дата обращения:03.01.2026

4. Christofferson S.A., McNish R.S., Sias D.L. Where mergers go
wrong/The McKinsey Quarterly. -- 2004. - № 2.- P.1-7.

5. Andrade G., Mitchell M., Stafford E. New Evidence and Perspectives on
Mergers// Journal of Economic Perspectives. -2001. -Vol.15 (2). -
P.103-120. \href{https://doi.org/10.1257/jep.15.2.103}{DOI
10.1257/jep.15.2.103}.

6. Andrade G., Stafford E. Investigating the economic role of mergers//
Journal of Corporate Finance. - 2004.- Vol.10(1). -P.1-36 DOI
\href{https://doi.org/10.1016/S0929-1199(02)00023-8}{10.1016/S0929-1199(02)00023-8}.

7. Breinlich H. Trade liberalization and industrial restructuring
through mergers and acquisitions// Journal of International Economics.
-2008. -Vol.76(2). - P.254-266. DOI
\href{https://doi.org/10.1016/j.jinteco.2008.07.007}{10.1016/j.jinteco.2008.07.007}.

8. UNCTAD World Investment Report. International investment in the
digital economy. -- 2025.-262 p. ISBN 978-92-1-003558-3.

9. MacKinlay A.C. Event Studies in Economics and Finance// Journal of
Economic Literature. -1997.-Vol.35(1). - P.13-39.

10. Moeller S.B., Schlingemann F.P., Stulz R.M. Firm size and the gains
from acquisitions// Journal of Financial Economics. -2004. -Vol.73(2).
-P.201--228. DOI
\href{https://doi.org/10.1016/j.jfineco.2003.07.002}{10.1016/j.jfineco.2003.07.002}.

11. Rau P.R., Vermaelen T. Glamour, value and the post-acquisition
performance of acquiring firms// Journal of Financial Economics. -1998.
-Vol.49(2). - P.223--253. DOI
\href{https://doi.org/10.1016/S0304-405X(98)00023-3}{10.1016/S0304-405X(98)00023-3}.

12. Wangerin D. M\&A Due Diligence, Post-Acquisition Performance, and
Financial Reporting for Business Combinations// Contemporary Accounting
Research. -2019.-Vol.36(4). - P.2344--2378. DOI
\href{https://doi.org/10.1111/1911-3846.12520}{10.1111/1911-3846.12520}.

13. Golubov A., Petmezas D., Travlos N.G. When It Pays to Pay Your
Investment Banker: New Evidence on the Role of Financial Advisors in
M\&As// The Journal of Finance. -2012. -Vol.67(1). P.271--311. DOI
\href{https://doi.org/10.1111/j.1540-6261.2011.01712.x}{10.1111/j.1540-6261.2011.01712.x}.

14. Rau P.R. Investment bank market share, contingent fee payments, and
the performance of acquiring firms// Journal of Financial Economics.
-2000. -Vol.56(2). - P.293--324.
DOI:\href{https://doi.org/10.1016/S0304-405X(00)00042-8}{10.1016/S0304-405X(00)00042-8}.

15. Potepa J., Thomas J. Goodwill Impairment After M\&A:
Acquisition-Level Evidence// SSRN Electronic Journal. -2022. -Vol.8(2).
-P.131-155. DOI
\href{http://doi.org/10.2139/ssrn.4239812}{10.2139/ssrn.4239812}.

16. Kothari S.P., Warner J.B. The Econometrics of Event Studies// SSRN
Electronic Journal. -2004. --Vol.1. DOI
\href{https://doi.org/10.2139/ssrn.608601}{10.2139/ssrn.608601}.

17. Neto J.Q.F., Bozos K., Dutordoir M., Nikolopoulos K. Are acquirer
stock price reactions to M\&A announcements in any way predictable? A
machine-learning analysis// Journal of the Operational Research Society.
-2025. -P.1-23. DOI
\href{https://doi.org/10.1080/01605682.2025.2562956}{10.1080/01605682.2025.2562956}.

\emph{{\bfseries Information about the author}}

Amanbekova A. - Master's student, Kazakh-British Technical University,
Almaty, Kazakhstan, e-mail:
amanbekovaidana@gmail.com.

\emph{{\bfseries Информация об авторе}}

Аманбекова А. М. - Магистрант, Казахстанско-Британский Технический
Университет, Алматы, Казахстан, e-mail:
amanbekovaidana@gmail.com.\